\lecture{2}{2026-09-11}{Recursive specifications and parameters}

When combinatorial classes are defined iteratively from the base classes \(\mathcal{E}\) and \(\mathcal{Z}\) using only the basic constructions (sum, product, and sequence), the resulting generating functions are always rational (i.e., ratios of two polynomials).

\section*{Recursive specifications}

When combinatorial classes are defined recursively in terms of themselves, their specifications yield algebraic or implicit equations for their generating functions.

\begin{example}[Rooted planar binary trees, {\cite[Example 1.2]{Melczer2021}, \cite[Chapter 5]{MelczerNotes}}]
    A \emph{rooted planar binary tree} is either empty or consists of a root node and two ordered binary subtrees (left and right).
    Let \(\mathcal{B}\) be the class of binary trees with size \(|\cdot|\) equal to the number of nodes.
    The recursive specification is
    \[
        \mathcal{B} = \mathcal{E} + \mathcal{Z} \times \mathcal{B} \times \mathcal{B},
    \]
    which translates into the quadratic equation
    \[
        B(z) = 1 + z B(z)^2.
    \]
    Solving for \(B(z)\) with the boundary condition \(B(0) = 1\) (since the empty tree has size \(0\)) yields
    \[
        B(z) = \frac{1 - \sqrt{1-4z}}{2z}.
    \]
    We can extract the coefficients via the generalized binomial theorem:
    \begin{align*}
        b_n &= [z^n] B(z) = [z^{n+1}] \frac{1 - \sqrt{1-4z}}{2} \\
        &= -\frac{1}{2} \binom{1/2}{n+1} (-4)^{n+1} = \frac{1}{n+1} \binom{2n}{n}.
    \end{align*}
\end{example}

\begin{example}[Quinary trees, {\cite[Chapter 6]{MelczerNotes}}]
    Quinary (or \(5\)-ary) trees satisfy the recursive specification
    \[
        \mathcal{Q} = \mathcal{E} + \mathcal{Z} \times \mathcal{Q}^5 \implies Q(z) = 1 + z Q(z)^5.
    \]
    Since this equation of degree \(5\) is not solvable in radicals, we use implicit methods such as \emph{LIFT} (the Lagrange Implicit Function Theorem).
\end{example}

\section*{Classes with parameters}

\begin{definition}[Parameter, {\cite[Chapter 7]{MelczerNotes}}]
A \emph{parameter} on a combinatorial class \((\mathcal{C}, |\cdot|)\) is a function \(\chi \colon \mathcal{C} \to \mathbb{N}\) assigning a non-negative integer to each object.
Unlike the weight function \(|\cdot|\), the number of objects of a given parameter value is not required to be finite.
\end{definition}

\begin{definition}[Bivariate generating function, {\cite[Chapter 7]{MelczerNotes}}]
The \emph{bivariate generating function} of a class with parameter \((\mathcal{C}, |\cdot|, \chi)\) is the formal power series
\[
    C(u, z) = \sum_{\sigma \in \mathcal{C}} u^{\chi(\sigma)} z^{|\sigma|} = \sum_{n \ge 0} \left( \sum_{k \ge 0} c_{n,k} u^k \right) z^n,
\]
where \(c_{n,k} = |\{\sigma \in \mathcal{C} : |\sigma| = n, \chi(\sigma) = k\}|\).
\end{definition}

Setting \(u=1\) yields \(C(1, z) = \sum_{n \ge 0} c_n z^n\), which recovers the univariate generating function of \(\mathcal{C}\), while setting \(u=0\) gives \(C(0, z) = \sum_{n \ge 0} c_{n,0} z^n\), the generating function for objects with parameter value \(0\).

\subsection*{Tracking classes}
To track a parameter in combinatorial constructions, we introduce a \emph{tracking class} \(\mu\) containing a single element of size \(0\) and parameter value \(1\), with bivariate generating function \(u\).

\begin{example}[Binary strings parameterized by number of zeroes, {\cite[Chapter 7]{MelczerNotes}}]
    Let \(\mathcal{B}\) be binary strings with size being length, and parameter \(\chi\) being the number of zeroes.
    Replacing the atom \(\mathcal{Z}_0\) with \(\mathcal{Z}_0 \times \mu\):
    \[
        \mathcal{B} = \operatorname{SEQ}(\mathcal{Z}_0 \times \mu + \mathcal{Z}_1) \implies B(u, z) = \frac{1}{1 - (uz + z)} = \frac{1}{1 - (1+u)z}.
    \]
    Expanding:
    \[
        B(u, z) = \sum_{n \ge 0} (1+u)^n z^n = \sum_{n \ge 0} \left( \sum_{k=0}^n \binom{n}{k} u^k \right) z^n.
    \]
    As a check, setting \(u = 1\) yields \(B(1, z) = \frac{1}{1-2z} = \sum_{n \ge 0} 2^n z^n\), recovering the univariate generating function for binary strings (since \(\sum_{k=0}^n \binom{n}{k} = 2^n\)), while setting \(u = 0\) gives \(B(0, z) = \frac{1}{1-z} = \sum_{n \ge 0} z^n\), the generating function for binary strings with no zeroes (i.e., strings consisting entirely of ones).
\end{example}

\begin{example}[Rooted planar binary trees parameterized by leaves, {\cite[Chapter 7]{MelczerNotes}}]
    Let \(\mathcal{B}\) be the class of binary trees with size equal to the number of nodes, and parameter \(\chi\) counting the number of leaves (nodes with two empty subtrees).
    
    The class of non-empty binary trees \(\mathcal{N} = \mathcal{B} \setminus \mathcal{E}\) satisfies
    \[
        \mathcal{N} = \mathcal{Z} \times \mu + 2(\mathcal{Z} \times \mathcal{N}) + \mathcal{Z} \times \mathcal{N} \times \mathcal{N},
    \]
    where:
    \begin{itemize}
        \item \(\mathcal{Z} \times \mu\): a leaf node (both children empty);
        \item \(2(\mathcal{Z} \times \mathcal{N})\): root with one empty child and one non-empty child (the root is not a leaf; leaves come from \(\mathcal{N}\));
        \item \(\mathcal{Z} \times \mathcal{N}^2\): root with two non-empty children (leaves come from both subtrees).
    \end{itemize}
    Its bivariate generating function satisfies
    \[
        N(u, z) = uz + 2z N(u, z) + z N(u, z)^2.
    \]
    Solving the quadratic equation \(z N^2 + (2z-1)N + uz = 0\) with \(N(u, 0) = 0\):
    \[
        N(u, z) = \frac{1 - 2z - \sqrt{1 - 4z + 4(1-u)z^2}}{2z}.
    \]
    Then \(B(u, z) = 1 + N(u, z)\).
\end{example}

\section*{Expected values and moments}

The \emph{expected value} of a parameter \(\chi\) over all objects of size \(n\) is (see \cite[Chapter 7]{MelczerNotes})
\[
    \mathbb{E}[\chi_n] = \sum_{k \ge 0} k \cdot \frac{c_{n,k}}{c_n} = \frac{1}{c_n} \sum_{k \ge 0} k c_{n,k}.
\]
Differentiating \(C(u, z)\) with respect to \(u\) and evaluating at \(u=1\):
\[
    C_u(1, z) = \left. \frac{\partial C(u, z)}{\partial u} \right|_{u=1} = \sum_{n \ge 0} \left( \sum_{k \ge 0} k c_{n,k} \right) z^n \implies \mathbb{E}[\chi_n] = \frac{[z^n] C_u(1, z)}{[z^n] C(1, z)}.
\]

\begin{example}[Expected number of leaves in a binary tree, {\cite[Chapter 7]{MelczerNotes}}]
    Differentiating \(N(u, z)\) at \(u=1\):
    \[
        N_u(1, z) = \left. \frac{\partial N(u, z)}{\partial u}\right|_{u=1} = \frac{z}{\sqrt{1-4z}} \implies [z^n] N_u(1, z) = \binom{2n-2}{n-1}.
    \]
    Thus, the expected number of leaves in a tree of size \(n\) is
    \[
        \mathbb{E}[\chi_n] = \frac{\binom{2n-2}{n-1}}{\frac{1}{n+1}\binom{2n}{n}} = \frac{n(n+1)}{2(2n-1)} \sim \frac{n}{4}.
    \]
\end{example}